\documentclass[12pt, a4paper]{article}

\usepackage{amsmath, amssymb, amsthm}
\usepackage{geometry}
\usepackage{fancyhdr}
\usepackage{enumitem}
\usepackage{xcolor}
\usepackage{hyperref}
\usepackage{tcolorbox}
\usepackage{xeCJK}
% ── 中文字体自动探测（由 render_latex.py 生成）──
% 探测到的候选字体：Source Han Serif SC, Noto Sans SC, SimSun, Microsoft YaHei, SimHei, KaiTi
\IfFontExistsTF{Source Han Serif SC}{\setCJKmainfont{Source Han Serif SC}}{}
\IfFontExistsTF{Noto Sans SC}{\setCJKmainfont{Noto Sans SC}}{}
\IfFontExistsTF{SimSun}{\setCJKmainfont{SimSun}}{}
\IfFontExistsTF{Microsoft YaHei}{\setCJKmainfont{Microsoft YaHei}}{}
\IfFontExistsTF{SimHei}{\setCJKmainfont{SimHei}}{}
\IfFontExistsTF{KaiTi}{\setCJKmainfont{KaiTi}}{}
\IfFontExistsTF{FangSong}{\setCJKmainfont{FangSong}}{}
\IfFontExistsTF{STSong}{\setCJKmainfont{STSong}}{}
% 若以上全部缺失，中文可能显示为方框；请安装 Noto Sans CJK SC 或在 config 中指定字体
\XeTeXlinebreaklocale "zh"
\XeTeXlinebreakskip = 0pt plus 1pt
% 关闭 CJK 与西文之间的额外间距，贴近学术排版习惯

\geometry{margin=1in}
\pagestyle{fancy}
\fancyhf{}
\lhead{MATH 221 Linear Algebra}
\rhead{Meteorain}
\rfoot{Page \thepage}
\cfoot{}

% ── 颜色定义 ──
\definecolor{solutionblue}{HTML}{1D4ED8}
\definecolor{solutiongreen}{HTML}{15803D}
\definecolor{warnamber}{HTML}{B45309}
\definecolor{verifyred}{HTML}{B91C1C}

% ── 自定义环境 ──
\newtcolorbox{stepbox}{colback=blue!2!white,colframe=blue!20!black,boxrule=0.6pt,arc=2pt,left=6pt,right=6pt,top=4pt,bottom=4pt}
\newtcolorbox{verifybox}{colback=green!3!white,colframe=green!35!black,boxrule=0.6pt,arc=2pt,left=6pt,right=6pt,top=3pt,bottom=3pt}
\newtcolorbox{warnbox}{colback=orange!5!white,colframe=orange!45!black,boxrule=0.6pt,arc=2pt,left=6pt,right=6pt,top=3pt,bottom=3pt}

\newcommand{\problem}[1]{\subsection*{Problem~#1}}
\newcommand{\solution}{%
  \par\noindent
  {\color{solutiongreen}\bfseries Solution.}\par
}
\newcommand{\verifyok}{\textcolor{solutiongreen}{\ensuremath{\checkmark}\ 已验证}}
\newcommand{\verifywarn}{\textcolor{verifyred}{\ensuremath{\times}\ 验证存疑}}

\begin{document}

\begin{center}
  {\LARGE\bfseries Homework 6 — Solutions} \\[0.6em]
  {\large\color{solutionblue} MATH 221 Linear Algebra} \\[0.4em]
  Meteorain \quad | \quad 2026-10-06
\end{center}
\vspace{0.5em}\hrule\vspace{1.2em}
\problem{1}

\begin{stepbox}
Compute the determinant of the matrix $A = \begin{pmatrix} 1 & 2 \\ 3 & 4 \end{pmatrix}$.
\end{stepbox}

\solution

\begin{stepbox}
解析到矩阵 $A = \left[\begin{matrix}1 & 2\\3 & 4\end{matrix}\right]$（阶数 2x2）
\end{stepbox}

\begin{stepbox}
行列式定义：按第一行展开 $\det(A) = -2$
\end{stepbox}

\begin{center}
$\boxed{-2}$
\end{center}

\begin{verifybox}\small \verifyok\quad 1/1 项检验通过（答案形式检查） \end{verifybox}

\vspace{0.8em}\hrule\vspace{0.8em}

\problem{2}

\begin{stepbox}
Let $A = \begin{pmatrix} 4 & 1 \\ 2 & 3 \end{pmatrix}$. Find $A^{-1}$.
\end{stepbox}

\solution

\begin{stepbox}
解析到矩阵 $A = \left[\begin{matrix}4 & 1\\2 & 3\end{matrix}\right]$（阶数 2x2）
\end{stepbox}

\begin{stepbox}
$\det(A) = 10$ → 非零，逆矩阵存在
\end{stepbox}

\begin{stepbox}
$A^{-1} = \left[\begin{matrix}\frac{3}{10} & - \frac{1}{10}\\- \frac{1}{5} & \frac{2}{5}\end{matrix}\right]$
\end{stepbox}

\begin{stepbox}
验证 $AA^{-1} = I$：$\begin{pmatrix}\left[\begin{matrix}1 & 0\\0 & 1\end{matrix}\right]\end{pmatrix}$
\end{stepbox}

\begin{center}
$\boxed{\left[\begin{matrix}\frac{3}{10} & - \frac{1}{10}\\- \frac{1}{5} & \frac{2}{5}\end{matrix}\right]}$
\end{center}

\begin{verifybox}\small \verifyok\quad 3/3 项检验通过（答案形式检查、矩阵形状检查、矩阵恒等式验证） \end{verifybox}

\vspace{0.8em}\hrule\vspace{0.8em}

\problem{3}

\begin{stepbox}
Find all eigenvalues of $A = \begin{pmatrix} 2 & 1 \\ 1 & 2 \end{pmatrix}$.
\end{stepbox}

\solution

\begin{stepbox}
解析到矩阵 $A = \left[\begin{matrix}2 & 1\\1 & 2\end{matrix}\right]$（阶数 2x2）
\end{stepbox}

\begin{stepbox}
特征多项式：$\chi_A(\lambda) = \lambda^{2} - 4 \lambda + 3$
\end{stepbox}

\begin{stepbox}
特征值：$\lambda = 3$，重数 1；$\lambda = 1$，重数 1
\end{stepbox}

\begin{center}
$\boxed{3, 1}$
\end{center}

\begin{verifybox}\small \verifyok\quad 2/2 项检验通过（答案形式检查、特征值特征多项式验证） \end{verifybox}

\vspace{0.8em}\hrule\vspace{0.8em}

\problem{4}

\begin{stepbox}
Find the eigenvalues and corresponding eigenvectors of
$A = \begin{pmatrix} 1 & 0 \\ 0 & 1 \end{pmatrix}$.
\end{stepbox}

\solution

\begin{stepbox}
解析到矩阵 $A = \left[\begin{matrix}1 & 0\\0 & 1\end{matrix}\right]$（阶数 2x2）
\end{stepbox}

\begin{stepbox}
$\lambda=1$ 的特征空间基：\textbackslash{}left[\textbackslash{}begin\{matrix\}1\textbackslash{}\textbackslash{}0\textbackslash{}end\{matrix\}\textbackslash{}right], \textbackslash{}left[\textbackslash{}begin\{matrix\}0\textbackslash{}\textbackslash{}1\textbackslash{}end\{matrix\}\textbackslash{}right]
\end{stepbox}

\begin{center}
$\boxed{\lambda = 1: \; \left[\begin{matrix}1\\0\end{matrix}\right], \left[\begin{matrix}0\\1\end{matrix}\right]}$
\end{center}

\begin{verifybox}\small \verifyok\quad 1/1 项检验通过（答案形式检查） \end{verifybox}

\vspace{0.8em}\hrule\vspace{0.8em}

\problem{5}

\begin{stepbox}
Compute the eigenvalues of $A = \begin{pmatrix} 1 & 2 & 0 \\ 0 & 1 & 0 \\ 0 & 0 & 5 \end{pmatrix}$.
\end{stepbox}

\solution

\begin{stepbox}
解析到矩阵 $A = \left[\begin{matrix}1 & 2 & 0\\0 & 1 & 0\\0 & 0 & 5\end{matrix}\right]$（阶数 3x3）
\end{stepbox}

\begin{stepbox}
特征多项式：$\chi_A(\lambda) = \lambda^{3} - 7 \lambda^{2} + 11 \lambda - 5$
\end{stepbox}

\begin{stepbox}
特征值：$\lambda = 1$，重数 2；$\lambda = 5$，重数 1
\end{stepbox}

\begin{center}
$\boxed{1\ (2\times), 5}$
\end{center}

\begin{verifybox}\small \verifyok\quad 2/2 项检验通过（答案形式检查、特征值特征多项式验证） \end{verifybox}

\vspace{0.8em}\hrule\vspace{0.8em}

\problem{6}

\begin{stepbox}
Let $A = \begin{pmatrix} 1 & 2 \\ 3 & 4 \end{pmatrix}$ and $\mathbf{b} = \begin{pmatrix} 5 \\ 11 \end{pmatrix}$.
Solve the linear system $A\mathbf{x} = \mathbf{b}$.
\end{stepbox}

\solution

\begin{stepbox}
解析到矩阵 $A = \left[\begin{matrix}1 & 2\\3 & 4\end{matrix}\right]$（阶数 2x2）
\end{stepbox}

\begin{stepbox}
$A = \left[\begin{matrix}1 & 2\\3 & 4\end{matrix}\right]$, $\mathbf{b} = \left[\begin{matrix}5\\11\end{matrix}\right]$
\end{stepbox}

\begin{stepbox}
$\det(A) = -2$ → 系数矩阵非奇异，唯一解
\end{stepbox}

\begin{stepbox}
$x = A^{-1}\mathbf{b} = \left[\begin{matrix}-2 & 1\\\frac{3}{2} & - \frac{1}{2}\end{matrix}\right] \cdot \left[\begin{matrix}5\\11\end{matrix}\right] = \left[\begin{matrix}1\\2\end{matrix}\right]$
\end{stepbox}

\begin{stepbox}
验证 $Ax - \mathbf{b} = \left[\begin{matrix}0\\0\end{matrix}\right]$
\end{stepbox}

\begin{center}
$\boxed{\left[\begin{matrix}1\\2\end{matrix}\right]}$
\end{center}

\begin{verifybox}\small \verifyok\quad 3/3 项检验通过（答案形式检查、矩阵形状检查、矩阵恒等式验证） \end{verifybox}

\vspace{0.8em}\hrule\vspace{0.8em}

\problem{7}

\begin{stepbox}
Determine whether the columns of
$A = \begin{pmatrix} 1 & 1 \\ 0 & 1 \end{pmatrix}$ are orthogonal.
\end{stepbox}

\solution

\begin{stepbox}
解析到矩阵 $A = \left[\begin{matrix}1 & 1\\0 & 1\end{matrix}\right]$（阶数 2x2）
\end{stepbox}

\begin{stepbox}
输入为 2x2 矩阵，取其列向量作为待处理向量：v\_1 = \textbackslash{}left[\textbackslash{}begin\{matrix\}1\textbackslash{}\textbackslash{}0\textbackslash{}end\{matrix\}\textbackslash{}right], v\_2 = \textbackslash{}left[\textbackslash{}begin\{matrix\}1\textbackslash{}\textbackslash{}1\textbackslash{}end\{matrix\}\textbackslash{}right]
\end{stepbox}

\begin{stepbox}
两两内积：$\langle v_1,v_2\rangle = 1$
\end{stepbox}

\begin{stepbox}
结论：**不正交**——存在内积非零的一对
\end{stepbox}

\begin{center}
**不正交**——存在内积非零的一对;\textbackslash{}quad \$\textbackslash{}langle v\_1,v\_2\textbackslash{}rangle = 1\$
\end{center}

\begin{verifybox}\small \verifyok\quad 1/1 项检验通过（答案形式检查） \end{verifybox}

\vspace{0.8em}\hrule\vspace{0.8em}

\problem{8}

\begin{stepbox}
Apply Gram-Schmidt orthogonalization to
$A = \begin{pmatrix} 1 & 1 \\ 1 & 0 \\ 0 & 1 \end{pmatrix}$.
\end{stepbox}

\solution

\begin{stepbox}
解析到矩阵 $A = \left[\begin{matrix}1 & 1\\1 & 0\\0 & 1\end{matrix}\right]$（阶数 3x2）
\end{stepbox}

\begin{stepbox}
输入为 3x2 矩阵，取其行向量作为待处理向量：v\_1 = \textbackslash{}left[\textbackslash{}begin\{matrix\}1 \& 1\textbackslash{}end\{matrix\}\textbackslash{}right], v\_2 = \textbackslash{}left[\textbackslash{}begin\{matrix\}1 \& 0\textbackslash{}end\{matrix\}\textbackslash{}right], v\_3 = \textbackslash{}left[\textbackslash{}begin\{matrix\}0 \& 1\textbackslash{}end\{matrix\}\textbackslash{}right]
\end{stepbox}

\begin{stepbox}
Gram-Schmidt 结果：
\end{stepbox}

\begin{stepbox}
$q_1 = \left[\begin{matrix}1\\1\\0\end{matrix}\right]$
\end{stepbox}

\begin{stepbox}
$q_2 = \left[\begin{matrix}\frac{1}{2}\\- \frac{1}{2}\\1\end{matrix}\right]$
\end{stepbox}

\begin{stepbox}
验证 $Q^TQ = \left[\begin{matrix}2 & 0\\0 & \frac{3}{2}\end{matrix}\right]$
\end{stepbox}

\begin{stepbox}
说明：$Q$ 的**各列**依次为 $q_1, q_2, \dots$，即所求正交组。
\end{stepbox}

\begin{center}
$\boxed{Q = \left[\begin{matrix}1 & \frac{1}{2}\\1 & - \frac{1}{2}\\0 & 1\end{matrix}\right]}$
\end{center}

\begin{verifybox}\small \verifyok\quad 2/2 项检验通过（答案形式检查、矩阵恒等式验证） \end{verifybox}

\vspace{0.8em}\hrule\vspace{0.8em}

\problem{9}

\begin{stepbox}
Compute the rank of
$A = \begin{pmatrix} 1 & 2 \\ 2 & 4 \\ 3 & 6 \end{pmatrix}$.
\end{stepbox}

\solution

\begin{stepbox}
解析到矩阵 $A = \left[\begin{matrix}1 & 2\\2 & 4\\3 & 6\end{matrix}\right]$（阶数 3x2）
\end{stepbox}

\begin{stepbox}
行阶梯化后非零行数为 1
\end{stepbox}

\begin{center}
$\boxed{1}$
\end{center}

\begin{verifybox}\small \verifyok\quad 1/1 项检验通过（答案形式检查） \end{verifybox}

\vspace{0.8em}\hrule\vspace{0.8em}

\problem{10}

\begin{stepbox}
Let $A = \begin{pmatrix} 2 & 0 \\ 0 & 3 \end{pmatrix}$. Find a matrix $P$ that diagonalizes
$A$, i.e. find $P$ and $D$ such that $A = PDP^{-1}$.
\end{stepbox}

\solution

\begin{stepbox}
解析到矩阵 $A = \left[\begin{matrix}2 & 0\\0 & 3\end{matrix}\right]$（阶数 2x2）
\end{stepbox}

\begin{stepbox}
$P = \left[\begin{matrix}1 & 0\\0 & 1\end{matrix}\right]$
\end{stepbox}

\begin{stepbox}
$\det(P) = 1$
\end{stepbox}

\begin{stepbox}
$D = \left[\begin{matrix}2 & 0\\0 & 3\end{matrix}\right]$
\end{stepbox}

\begin{stepbox}
验证 $P^{-1}AP = \left[\begin{matrix}2 & 0\\0 & 3\end{matrix}\right]$
\end{stepbox}

\begin{stepbox}
✓ 与 D 一致，对角化成立
\end{stepbox}

\begin{center}
$\boxed{A = PDP^{-1},\quad P = \left[\begin{matrix}1 & 0\\0 & 1\end{matrix}\right],\quad D = \left[\begin{matrix}2 & 0\\0 & 3\end{matrix}\right]}$
\end{center}

\begin{verifybox}\small \verifyok\quad 4/4 项检验通过（答案形式检查、矩阵形状检查、矩阵恒等式验证、特征值特征多项式验证） \end{verifybox}

\vspace{0.8em}\hrule\vspace{0.8em}

\problem{11}

\begin{stepbox}
Let $A = \begin{pmatrix} 0 & 0 \\ 0 & 0 \end{pmatrix}$. What is the null space of $A$?
\end{stepbox}

\solution

\begin{stepbox}
解析到矩阵 $A = \left[\begin{matrix}0 & 0\\0 & 0\end{matrix}\right]$（阶数 2x2）
\end{stepbox}

\begin{stepbox}
基：\textbackslash{}left[\textbackslash{}begin\{matrix\}1\textbackslash{}\textbackslash{}0\textbackslash{}end\{matrix\}\textbackslash{}right], \textbackslash{}left[\textbackslash{}begin\{matrix\}0\textbackslash{}\textbackslash{}1\textbackslash{}end\{matrix\}\textbackslash{}right]
\end{stepbox}

\begin{center}
$\boxed{\left\{ \left[\begin{matrix}1\\0\end{matrix}\right], \left[\begin{matrix}0\\1\end{matrix}\right] \right\}}$
\end{center}

\begin{verifybox}\small \verifyok\quad 1/1 项检验通过（答案形式检查） \end{verifybox}

\vspace{0.8em}\hrule\vspace{0.8em}

\problem{12}

\begin{stepbox}
Diagonalize $A = \begin{pmatrix} 4 & 1 \\ 2 & 3 \end{pmatrix}$.
\end{stepbox}

\solution

\begin{stepbox}
解析到矩阵 $A = \left[\begin{matrix}4 & 1\\2 & 3\end{matrix}\right]$（阶数 2x2）
\end{stepbox}

\begin{stepbox}
$P = \left[\begin{matrix}-1 & 1\\2 & 1\end{matrix}\right]$
\end{stepbox}

\begin{stepbox}
$\det(P) = -3$
\end{stepbox}

\begin{stepbox}
$D = \left[\begin{matrix}2 & 0\\0 & 5\end{matrix}\right]$
\end{stepbox}

\begin{stepbox}
验证 $P^{-1}AP = \left[\begin{matrix}2 & 0\\0 & 5\end{matrix}\right]$
\end{stepbox}

\begin{stepbox}
✓ 与 D 一致，对角化成立
\end{stepbox}

\begin{center}
$\boxed{A = PDP^{-1},\quad P = \left[\begin{matrix}-1 & 1\\2 & 1\end{matrix}\right],\quad D = \left[\begin{matrix}2 & 0\\0 & 5\end{matrix}\right]}$
\end{center}

\begin{verifybox}\small \verifyok\quad 4/4 项检验通过（答案形式检查、矩阵形状检查、矩阵恒等式验证、特征值特征多项式验证） \end{verifybox}

\vspace{0.8em}\hrule\vspace{0.8em}

\vspace{1.5em}
\hrule\vspace{0.8em}
\begin{center}\small
\textbf{求解汇总}\quad 共 12 题，自动解出 12 题（100\%）；验证通过 12 题，存疑 0 题。\\[0.3em]
\textcolor{gray}{求解器分布：sympy\_diagonalize×2、sympy\_eigenvalues×2、sympy\_eigenvectors×1、sympy\_gram\_schmidt×1、sympy\_linear\_system×1、sympy\_matrix\_determinant×1、sympy\_matrix\_inverse×1、sympy\_matrix\_rank×1、sympy\_nullspace×1、sympy\_orthogonality\_check×1}
\end{center}
\vfill\begin{center}\footnotesize\textcolor{gray}{由 Meteorain C4C 作业自动求解流水线自动生成}\end{center}
\end{document}
